% ---------------------------------------------------------------------------
% Ch. 8 §2 -- Selection of the DER Q(V) dead-zone half-width.
% METHOD SECTION. Results are placeholders (marked RESULTS PLACEHOLDER).
%
% To \input this into the thesis, delete everything from \documentclass down to
% \tableofcontents, and the closing \end{document}. The macros in the preamble
% (\pu, \Mvar, \qss, \rms) match docs/qss_rms_divergence_analysis.tex; drop them
% if the thesis already defines equivalents.
%
% Generated 2026-07-31 alongside the runs it describes.
% ---------------------------------------------------------------------------
\documentclass[11pt,a4paper]{article}

\usepackage[utf8]{inputenc}
\usepackage[T1]{fontenc}
\usepackage{amsmath,amssymb}
\usepackage{booktabs}
\usepackage{siunitx}
\usepackage{graphicx}
\usepackage[margin=2.5cm]{geometry}
\usepackage[hidelinks]{hyperref}

\sisetup{detect-weight=true,detect-family=true}

\newcommand{\pu}{\ensuremath{\mathrm{p.u.}}}
\newcommand{\Mvar}{\ensuremath{\mathrm{Mvar}}}
\newcommand{\qss}{\textsc{qss}}
\newcommand{\rms}{\textsc{rms}}
\newcommand{\dz}{\ensuremath{\delta}}

\title{Selection of the DER $Q(V)$ dead-zone half-width\\[2mm]
\large Experimental design, choice of profile windows, and their
representativeness}
\author{}
\date{2026--07--31}

\begin{document}
\maketitle
\tableofcontents

%======================================================================
\section{Purpose and claim under test}

The distribution-connected DER follow a local reactive-power droop
$Q(V)$ with a dead zone of half-width $\dz$ about a re-anchored voltage
reference. The dead zone is a design parameter, and the claim to be tested is
that it is a \emph{two-sided} one:

\begin{description}
  \item[$\dz$ too narrow] the local droop answers every small profile
    variation. The DER move continuously and the outer online feedback
    optimisation (OFO) layer repeatedly re-anchors against them.
  \item[$\dz$ too wide] the droop is inactive between dispatches, so no local
    voltage support acts and the interface reactive power drifts until the next
    OFO step corrects it.
\end{description}

If both extremes degrade the controlled quantities, the useful $\dz$ is an
interior minimum, and the experiment must exhibit that minimum rather than
assume it.

\subsection{Assumptions, actuators and controlled outputs}

\begin{itemize}
  \item \textbf{Plant.} Modified IEEE~39-bus transmission system, scenario
    \texttt{rural\_700}, with four HV (\SI{110}{\kilo\volt}) distribution
    underlays. $\mathrm{DSO}_3$ carries a scenario multiplier of $2\times$
    installed DER and $2\times$ active-load base; its reactive load is
    \emph{not} scaled. Every DSO has zero constant reactive load and follows
    $Q_\mathrm{load}(t) = \SI{500}{\Mvar}\cdot m_\mathrm{qload}(t)$.
  \item \textbf{Actuators.} DER reactive power (continuous, $Q(V)$ droop with
    dead zone), coupler three-winding OLTCs, machine two-winding OLTCs, and
    MSC/MSR switched shunts.
  \item \textbf{Controlled outputs.} Reactive power flow at the EHV--HV
    interfaces, and nodal voltages within their admissible band.
  \item \textbf{DER capability.} The physical VDE-AR-N-4120 operating diagrams
    are used throughout (\texttt{--physical-capability}). This matters for the
    window selection: the diagram makes reactive capability contingent on
    active power, so a park at $P=0$ has \emph{no} reactive capability. See
    \S\ref{sec:degenerate}.
  \item \textbf{Disturbance class.} Profiled operation --- the annual
    load/DER time series, not an injected step. This is deliberate: the dead
    zone exists to filter the small-signal variation of normal operation, so a
    step test would not exercise the mechanism under study.
  \item \textbf{No contingencies, and no measurement noise.} Every run in the
    study records \texttt{contingencies = []},
    \texttt{measurement\_noise.enabled = False} and
    \texttt{enable\_reachability\_guard = False}. The profile evolution is
    therefore the \emph{only} exogenous disturbance. The consequences for the
    parameter under study are discussed in \S\ref{sec:excitation}; they are not
    incidental.
  \item \textbf{Controllers see only their cached sensitivities.} No controller
    is given the plant equations.
\end{itemize}

\subsection{Measured quantities}

Three controlled quantities are evaluated on the \emph{same} runs:

\begin{table}[htbp]
\centering
\begin{tabular}{lll}
\toprule
Symbol & Definition & Unit \\
\midrule
interface $Q$ & mean $|Q_\mathrm{act}-Q_\mathrm{set}|$ over the TS--DSO
  interfaces & \si{\Mvar} \\
TS voltage & per-zone RMS voltage error & \pu \\
DS voltage & RMS deviation of each DSO group's mean $V$ from
  \SI{1.03}{\pu} & \pu \\
\bottomrule
\end{tabular}
\caption{Controlled quantities evaluated for every dead band.}
\label{tab:metrics}
\end{table}

%======================================================================
\section{Experimental grid}

The sweep is a full factorial over dead band and profile window,
\[
  \dz \in \{0,\ 0.001,\ 0.0025,\ 0.005,\ 0.0075,\ 0.01,\ 0.015,\ 0.02,\ 0.03\}
  \ \pu,
\]
crossed with the profile windows of Table~\ref{tab:windows}. Each cell is an
independent closed-loop co-simulation against the PowerFactory \rms{} plant,
with a quasi-static reference leg run on the same controller configuration.

\paragraph{Why $\dz = 0.001$ was added, and what it showed.} The initial grid
stepped from $0$ straight to $0.0025$, across which the interface-$Q$ error
changes by a factor of three (\S\ref{sec:findings}). A single interval carrying
the entire narrow-side transition cannot distinguish a gradual degradation from
an abrupt one, and the distinction matters for the design rule: an abrupt
collapse means the dead band must simply be non-zero, whereas a gradual one
means its exact value trades off continuously against voltage quality.

The transition is \textbf{abrupt}. In the mildest window the interface-$Q$ error
falls from \SI{1.395}{\Mvar} at $\dz=0$ to \SI{0.533}{\Mvar} at
$\dz=\SI{0.001}{\pu}$ --- a dead zone of one part in a thousand recovers
\SI{90}{\percent} of the entire penalty between $\dz=0$ and the optimum
\SI{0.441}{\Mvar} at $\dz=0.005$. The remaining \SI{10}{\percent} is spread over
a range four times wider.

This changes the form of the design rule. The requirement is not a particular
$\dz$ but a \emph{non-zero} one: the tracking benefit saturates almost
immediately, while the DS voltage cost continues to accrue across the whole
range (\SI{0.00225}{\pu} at $\dz=0$, \SI{0.00418}{\pu} at $0.001$,
\SI{0.00902}{\pu} at $0.03$). A dead band should therefore be chosen as small as
the measurement chain permits rather than optimised for tracking --- with the
caveat of \S\ref{sec:excitation} that no measurement noise is present in these
runs, which is precisely what would set that lower limit in practice.

The grid was not swept in one pass. The initial matrix covered
$\dz\in\{0.0025,\dots,0.015\}$; $\dz=0$ and the upper points
$\dz\in\{0.02,0.03\}$ were added afterwards, the latter because the argument
minimum for one window was found to lie \emph{at} the upper end of the original
range and was therefore a bound rather than a measured optimum
(\S\ref{sec:edge}).

A note on the $\dz=0$ cell, since it decides the central claim. It was
initially excluded on the basis of a recorded cost of approximately
\SI{4.9}{\hour} per run, attributed to the static-leg initialisation at a zero
dead band. \textbf{That figure is void}: measured on 2026-07-31, the three
$\dz=0$ runs took \SI{12.3}{\minute}, \SI{12.8}{\minute} and
\SI{13.2}{\minute} --- indistinguishable from every other cell, and a factor of
roughly $23$ below the recorded estimate. The exclusion therefore withheld the
single most informative point in the study for no reason. All cells are
ordinary grid points and cost approximately \SI{13}{\minute}.

%======================================================================
\section{Selection of the profile windows}
\label{sec:windows}

\subsection{What the windows are for}

The dead band is a single scalar applied to all DER, so the question is not
only \emph{what} its optimum is but whether that optimum is a property of the
system or an artefact of the operating point at which it was measured. The
windows exist to answer that second question, and the design requirement is
therefore a \emph{spread} of operating points rather than a statistically
representative sample of the year.

\subsection{Provenance of the initial screening, and its limits}
\label{sec:screening}

The first three windows were drawn from a Tier-1 seasonal screening that ranked
quarter-hours by the quasi-static per-interval voltage excursion:

\begin{center}
\begin{tabular}{lll}
\toprule
Window & Screening excursion & Intended role \\
\midrule
2016-01-05 08:00 & \SI{0.00828}{\pu} & low (reference) \\
2016-01-15 03:00 & \SI{0.01573}{\pu} & mid \\
2016-07-15 03:00 & \SI{0.02051}{\pu} & high / annual maximum \\
\bottomrule
\end{tabular}
\end{center}

\textbf{These figures were measured on an earlier network topology and are not
valid descriptors of the runs reported here.} They are recorded because they
document how the windows were \emph{chosen}; they must not be used to
characterise the operating points, to order them by severity, or as a regressor
against which the optimum is fitted. \S\ref{sec:degenerate} gives the concrete
demonstration of why: the window this screening ranked as the annual maximum
turned out, on the current topology, to be the least demanding of the three.

The corresponding column in the machine-readable output is named
\texttt{screening\_excursion\_pu\_old\_topology} for this reason.

\subsection{Measured characterisation of the windows}

Table~\ref{tab:windows} characterises each window by quantities measured on the
current topology, either from the profile data or from the runs themselves.
These, not the screening figures, are what the windows should be described by.

\begin{table}[htbp]
\centering
\sisetup{table-number-alignment=center}
\begin{tabular}{lS[table-format=4.1]S[table-format=4.1]S[table-format=+5.1]cc}
\toprule
Window & {DER $P$} & {Load $P$} & {Net infeed} & Parks with & PCC capability \\
       & {[\si{\mega\watt}]} & {[\si{\mega\watt}]} & {[\si{\mega\watt}]}
       & zero $Q$ cap. & {[\si{\Mvar}]} \\
\midrule
2016-02-22 13:00 & 2262.4 & 2379.9 & -117.4 & {pending}
  & {pending} \\
2016-01-05 08:00 & 2324.5 & 1916.0 &  408.6 & 16/44
  & $[-45.2,\,41.6]$ \\
2016-01-15 03:00 & 2605.7 & 1800.7 &  805.0 & 16/44
  & $[-44.3,\,41.8]$ \\
2016-12-18 14:00 & 2741.8 & 1374.7 & 1367.1 & 16/44
  & $[-42.0,\,46.1]$ \\
\addlinespace
2016-07-15 03:00 &   29.2 & 1055.3 & -1026.1 & \textbf{44/44}
  & $[0.0,\,0.0]$ \\
2016-05-01 16:00 & 3458.6 & 1258.6 & 2200.0 & \multicolumn{2}{c}{\emph{no
  convergent solution}} \\
\bottomrule
\end{tabular}
\caption{Windows characterised by current-topology quantities, ordered by net
infeed. The first four are the live windows of the study. The last two carry no
results: 2016-07-15 03:00 is retained as a null result
(\S\ref{sec:degenerate}), and 2016-05-01 16:00 admits no convergent power-flow
solution (\S\ref{sec:extra-windows}). PCC capability is the initial zone-2
value; zone~3 is comparable. Net infeed is
$\sum P_\mathrm{DER}-\sum P_\mathrm{load}$ over the profiled rows.}
\label{tab:windows}
\end{table}

The three original live windows span a factor of $3.3$ in net infeed ($+409$ to
$+\SI{1367}{\mega\watt}$) while holding DER reactive capability essentially
constant: all three have the same $16/44$ parks below the VDE capability
threshold and comparable PCC envelopes of roughly $\pm\SI{45}{\Mvar}$. The
comparison therefore varies the loading condition while leaving the available
actuator authority unchanged, which is the property that makes a shift in
$\dz^\star$ attributable to the operating point rather than to a change in the
actuator set.

It should be stated precisely \emph{what} varies. Across those three, DER infeed
is similar ($2325$, $2606$, $\SI{2742}{\mega\watt}$) and the spread in net
infeed comes predominantly from the load, which falls from $1916$ to
$\SI{1375}{\mega\watt}$. The axis is therefore better described as
\emph{decreasing load at near-constant generation} than as increasing
generation, and the resulting stress is a light-load voltage-rise condition.

\subsection{Two further windows, and one that does not exist}
\label{sec:extra-windows}

All three original live windows are net \emph{exporting}. The dead band was
therefore selected entirely within the voltage-rise regime, in which the local
droop absorbs reactive power. Two windows were added to widen that basis.

\paragraph{2016-05-01 16:00 --- $+\SI{2200}{\mega\watt}$: attempted and
infeasible.} This window was chosen to extend the export axis
\SI{61}{\percent} above the previous maximum, deliberately below the annual
peak, because 2016-07-25 13:00 at $+\SI{3259}{\mega\watt}$ had already aborted
with a non-convergent Jacobian probe. \textbf{It failed the same way.} All nine
dead bands raised \texttt{LoadflowNotConverged} in the static leg, each within
roughly \SI{30}{\second} and before the simulation proper began, so the failure
is a property of the operating point and not of $\dz$.

The export axis therefore could not be extended, and the study has \emph{four}
live windows rather than five. What the attempt does buy is a much tighter
bound on feasibility: $+\SI{1367}{\mega\watt}$ converges and
$+\SI{2200}{\mega\watt}$ does not, where previously the bracket was
$+805$ to $+\SI{3259}{\mega\watt}$. A further attempt would sensibly target
$+1600$ to $+\SI{1700}{\mega\watt}$.

That this could not be predicted offline is itself worth recording: a
pandapower probe with a flat start declared feasible hours infeasible, while a
probe with runner-like settings declared this window feasible. Neither
reproduces the runner, so feasibility here is an empirical property established
by running.

\paragraph{2016-02-22 13:00 --- $-\SI{117}{\mega\watt}$, net import.} The only
window in the set with negative net infeed \emph{and} capable DER, and thus the
only one probing the voltage-drop regime, in which the droop must inject rather
than absorb.

\paragraph{A deep-import window with capable DER does not exist in this data.}
This is a property of the profiles rather than a gap in the search, and it
bounds what the study can claim. DER output and load are positively correlated
--- both peak in daytime hours --- so the hours with substantial reactive
capability are precisely the hours with high generation. Screening all
$35\,136$ quarter-hours, \emph{no} hour combines $P_\mathrm{DER} \geq
\SI{2200}{\mega\watt}$ with net infeed below $-\SI{300}{\mega\watt}$. The deeply
importing hours are the low-generation ones, and under the VDE capability
diagram their parks have little or no reactive capability --- the limiting case
being the degenerate window of \S\ref{sec:degenerate}, where infeed is
\SI{29}{\mega\watt} and capability is exactly zero.

The consequence is that \textbf{$\dz$ cannot be assessed under deep import with
active DER on this system}, because that operating point does not occur. Any
statement about dead-band selection in a voltage-drop regime is therefore
restricted to the near-balanced condition of 2016-02-22 13:00. A study of that
regime would need either a different generation mix or an artificial
dispatch, and either choice would leave the annual profile set.

\subsection{A degenerate window, and what it establishes}
\label{sec:degenerate}

The window originally selected as the most severe, 2016-07-15 03:00, returned
\emph{bit-identical} results for all five dead bands of the initial grid. The
identity extends to the full per-park reactive-power record, not merely to the
aggregated metrics, so it is not a numerical coincidence.

The cause is the DER capability model. At a July night hour there is no
photovoltaic infeed and the wind profiles sit inside a genuine multi-day lull,
giving an aggregate DER infeed of \SI{29.2}{\mega\watt} against
\SI{2605.7}{\mega\watt} at 2016-01-15 03:00. Under the VDE-AR-N-4120 diagram
reactive capability is contingent on active power, so \emph{all} 44 parks report
zero reactive capability and both zone PCC capability intervals collapse to
$[0,0]\,\si{\Mvar}$. With no DER able to move, the $Q(V)$ characteristic never
binds and its dead zone cannot influence anything.

Two conclusions follow, and both belong in the thesis rather than in a
footnote:

\begin{enumerate}
  \item \textbf{Dead-band selection presupposes DER reactive headroom.} The
    parameter is meaningful only where the continuous actuators it governs can
    act. This is a statement about the applicability of the design rule, not a
    limitation of the experiment.
  \item \textbf{The screening that selected this window was not merely
    imprecise but misleading}, having ranked it the most severe operating point
    of the three. This is the evidence for the provenance restriction stated in
    \S\ref{sec:screening}.
\end{enumerate}

The window is retained in the data set as a documented null result and is
excluded from the cross-window comparison of optima.

\subsection{Selection of the replacement window}

The replacement, 2016-12-18 14:00, was selected under three criteria applied in
order:

\begin{enumerate}
  \item \textbf{Actuator availability.} Every DER park must carry meaningful
    active power, so that reactive capability exists under the VDE diagram.
    Screening all $35\,136$ quarter-hours, this holds for approximately
    \SI{37}{\percent} of the year.
  \item \textbf{Elevated stress.} Net infeed must exceed that of the existing
    windows, net infeed being the quantity that drives the voltage rise the
    dead zone is meant to filter.
  \item \textbf{Numerical feasibility.} The operating point must admit a
    converged power-flow solution in the static leg.
\end{enumerate}

The third criterion is not a formality. A first candidate, 2016-07-25 13:00,
was selected on criteria 1--2 at $+\SI{3259}{\mega\watt}$ net infeed --- close to
the annual maximum --- and failed criterion 3: the static leg's Jacobian probe
did not converge. An offline power-flow screen was not able to predict this
reliably, since its verdict depended strongly on the initialisation and slack
treatment and did not reproduce the runner's behaviour. The criterion was
therefore applied empirically, and the replacement was chosen with margin
rather than at maximum stress: $+\SI{1367}{\mega\watt}$, a factor $1.7$ above
2016-01-15 03:00 and well inside the region known to solve.

\subsection{In what sense the windows are representative}
\label{sec:representative}

The claim that should be made is narrow, and it is worth stating precisely
because a stronger one is not supported.

\textbf{What the window set supports.} The four live windows are
\emph{contrasting} operating points spanning net infeed from
$-117$ to $+\SI{1367}{\mega\watt}$ (Table~\ref{tab:windows}), a range that
crosses zero and therefore covers both the voltage-rise and the near-balanced
condition. If the optimum $\dz^\star$ is common to all of them, that is evidence
of a system property; if it moves monotonically with loading, that is evidence
of an operating-point dependence with an identifiable direction. Either outcome
is informative, and both are conclusions about \emph{sensitivity to the
operating point}, which is what the design requires.

\textbf{What it does not support.} The windows are not a random or stratified
sample of the operating year, and no frequency-weighted statement --- such as an
expected annual performance, or a $\dz$ optimal in an annual-average sense ---
can be drawn from four purposively chosen points. Three of them were selected
under a screening now known to be invalid as a severity measure
(\S\ref{sec:screening}), and the fourth under a different criterion (net
infeed). The set is \emph{purposive}, not representative in the sampling sense,
and is reported as such. It is also asymmetric about zero: the export side
reaches $+\SI{1367}{\mega\watt}$ while the import side stops at
$-\SI{117}{\mega\watt}$. Both limits are structural rather than chosen --- the
import side by the correlation between DER output and load
(\S\ref{sec:extra-windows}), the export side by power-flow feasibility.

\textbf{What would strengthen it.} A screening of the annual profile on the
current topology, using a quantity measured on these runs --- DER reactive
headroom, or a per-interval voltage excursion recomputed post-fix --- would
allow the windows to be placed on a defensible severity axis and would permit a
stratified selection. That screening does not presently exist; the old one is
void, and inventing a value for the replacement window would mix two
incompatible measurements. This is recorded as an open item rather than
papered over.

%======================================================================
\section{Results}

% ===================== RESULTS PLACEHOLDER ============================
% Populate from results/deadband_selection/ once the full grid is complete:
%   deadband_metrics.csv  -- one row per run
%   deadband_optima.csv   -- argmin per window per metric, with at_range_edge
%   figures/deadband_selection.{png,pdf}
%   figures/deadband_band.{png,pdf}
% ======================================================================

\subsection{Per-window response}

\begin{figure}[htbp]
\centering
% \includegraphics[width=\textwidth]{figures/deadband_selection.pdf}
\framebox[0.9\textwidth][c]{\rule{0pt}{6cm}%
  \textsf{PLACEHOLDER: deadband\_selection.pdf}}
\caption{The three controlled quantities against the dead-zone half-width
$\dz$, one curve set per profile window. Left axis \si{\Mvar}, right axis
\pu. Each interior minimum is ringed.}
\label{fig:selection}
\end{figure}

\begin{table}[htbp]
\centering
\sisetup{table-number-alignment=center}
\begin{tabular}{S[table-format=1.4]
                S[table-format=1.3]S[table-format=1.5]S[table-format=1.5]
                S[table-format=1.3]S[table-format=1.5]S[table-format=1.5]
                S[table-format=1.3]S[table-format=1.5]S[table-format=1.5]}
\toprule
& \multicolumn{3}{c}{2016-01-05 08:00}
& \multicolumn{3}{c}{2016-01-15 03:00}
& \multicolumn{3}{c}{2016-12-18 14:00} \\
\cmidrule(lr){2-4}\cmidrule(lr){5-7}\cmidrule(lr){8-10}
{$\dz$} & {$Q$} & {TS $V$} & {DS $V$}
        & {$Q$} & {TS $V$} & {DS $V$}
        & {$Q$} & {TS $V$} & {DS $V$} \\
\midrule
0      & 1.395 & 0.00559 & 0.00225 & 4.367 & 0.00662 & 0.00821 & 3.568 & 0.00809 & 0.00577 \\
0.0025 & 0.463 & 0.00530 & 0.00535 & 2.824 & 0.00593 & 0.01199 & 1.838 & 0.00737 & 0.00991 \\
0.005  & 0.441 & 0.00531 & 0.00639 & 2.363 & 0.00586 & 0.01350 & 1.615 & 0.00726 & 0.01197 \\
0.0075 & 0.466 & 0.00549 & 0.00728 & 2.082 & 0.00589 & 0.01473 & 1.544 & 0.00736 & 0.01343 \\
0.01   & 0.510 & 0.00581 & 0.00776 & 1.951 & 0.00596 & 0.01587 & 1.512 & 0.00753 & 0.01419 \\
0.015  & 0.533 & 0.00619 & 0.00839 & 1.644 & 0.00601 & 0.01691 & 1.494 & 0.00749 & 0.01347 \\
0.02   & 0.527 & 0.00614 & 0.00887 & 1.524 & 0.00592 & 0.01650 & 1.210 & 0.00692 & 0.01377 \\
0.03   & 0.640 & 0.00609 & 0.00902 & 1.572 & 0.00582 & 0.01615 & 1.126 & 0.00703 & 0.01393 \\
\bottomrule
\end{tabular}
\caption{Measured response for every cell of the grid: interface $Q$ in
\si{\Mvar}, TS and DS voltage in \pu. The degenerate window
2016-07-15 03:00 is omitted --- all five of its cells returned
$Q = \SI{2.167}{\Mvar}$, TS $V = \SI{0.00592}{\pu}$,
DS $V = \SI{0.01570}{\pu}$ identically (\S\ref{sec:degenerate}).}
\label{tab:results}
\end{table}

\subsection{Does the optimum move with the operating point?}

\begin{figure}[htbp]
\centering
% \includegraphics[width=\textwidth]{figures/deadband_band.pdf}
\framebox[0.9\textwidth][c]{\rule{0pt}{6cm}%
  \textsf{PLACEHOLDER: deadband\_band.pdf}}
\caption{Minimum/maximum band across the live windows, each metric normalised
to its own best value, so that the \emph{shape} of the response can be compared
across operating points whose absolute levels differ.}
\label{fig:band}
\end{figure}

\begin{table}[htbp]
\centering
\begin{tabular}{lccc}
\toprule
Window & $\dz^\star$ (interface $Q$) & $\dz^\star$ (TS $V$) & $\dz^\star$ (DS $V$) \\
\midrule
2016-01-05 08:00 & 0.005 & 0.0025 & $0^{\dagger}$ \\
2016-01-15 03:00 & 0.02  & $0.03^{\dagger}$ & $0^{\dagger}$ \\
2016-12-18 14:00 & $0.03^{\dagger}$ & 0.02 & $0^{\dagger}$ \\
\addlinespace
2016-07-15 03:00 & \multicolumn{3}{c}{flat --- no optimum exists} \\
\bottomrule
\end{tabular}
\caption{Argument minima per window and metric, in \pu. $^{\dagger}$ marks an
entry at the edge of the swept range: a bound, not a measured optimum
(\S\ref{sec:edge}). The interface-$Q$ optima of the first two windows are
interior and therefore measured.}
\label{tab:optima}
\end{table}

\subsection{Findings}
\label{sec:findings}

\paragraph{The dead zone is genuinely two-sided, and the zero anchor is what
establishes it.} At $\dz=0$ the interface-$Q$ error is
$3.2\times$, $2.9\times$ and $3.2\times$ its optimal value in the three live
windows. Without that point the narrow-side branch rested on a difference of
about \SI{5}{\percent} between $\dz=0.0025$ and $\dz=0.005$ in the mildest
window, which is not separable from run-to-run variation. With it, the
narrow-side penalty is unambiguous.

\paragraph{$\dz^\star$ is strongly operating-point dependent.} The
interface-$Q$ optimum moves $0.005 \rightarrow 0.02 \rightarrow \geq 0.03$, a
factor of at least six. The first two are bracketed interior minima; the third
is still falling at the widest dead band swept, so the true separation is
larger than measured.

\paragraph{The two controlled quantities select opposite ends.} DS voltage is
best at $\dz=0$ in \emph{every} live window, while interface $Q$ is at its worst
there. Moving from $\dz=0$ to the interface-$Q$ optimum costs a factor
$2.0$--$2.8$ in DS voltage error. The dead zone therefore does not have a single
optimum across the controlled outputs: it purchases interface-$Q$ tracking at
the cost of DS voltage quality, and the design choice is where on that exchange
to sit. TS voltage does not discriminate --- its optimum lands at $0.0025$,
$0.03$ and $0.02$ across the three windows with a spread of \SIrange{1}{3}
{\percent}, and should not be used to select $\dz$.

\subsection{Pareto analysis: which dead bands survive both objectives}
\label{sec:pareto}

Since the two controlled quantities select opposite ends of the range, a single
$\dz^\star$ does not exist and the correct object is the set of non-dominated
dead bands. A dead band is \emph{dominated} when another is at least as good in
every objective and strictly better in one; a dominated dead band is never a
rational choice, because the alternative costs nothing to adopt.

\begin{table}[htbp]
\centering
\begin{tabular}{llll}
\toprule
Window & Net infeed & Pareto-optimal $\dz$ & dominated \\
\midrule
2016-02-22 13:00 & $-117$ & $0.005,\ 0.0075,\ 0.03$
  & $0,\ 0.001,\ 0.0025,\ 0.01,\ 0.015,\ 0.02$ \\
2016-01-05 08:00 & $+409$ & $0,\ 0.001,\ 0.0025,\ 0.005$
  & $0.0075,\ 0.01,\ 0.015,\ 0.02,\ 0.03$ \\
2016-01-15 03:00 & $+805$ & all but $0.015$ & $0.015$ \\
2016-12-18 14:00 & $+1367$ & all but $0.01$ & $0.01$ \\
\midrule
\textbf{intersection} & & $\mathbf{0.005}$ & \\
\bottomrule
\end{tabular}
\caption{Non-dominated dead bands per window over (interface $Q$, DS $V$), in
\pu; net infeed in \si{\mega\watt}. Including TS voltage as a third objective
leaves every set unchanged, which is independent confirmation that it does not
discriminate.}
\label{tab:pareto}
\end{table}

\paragraph{One dead band survives every regime.} The intersection is the single
value $\dz = \SI{0.005}{\pu}$. This is a considerably stronger statement than
the three export windows alone support: among those, $\dz \in \{0, 0.001,
0.0025, 0.005\}$ are all non-dominated, and no unique choice emerges. Adding the
one window with negative net infeed collapses the set to a point.

\paragraph{Why the import window is decisive: the voltage trade-off reverses.}
In every exporting window, DS voltage is monotonically best at $\dz = 0$ ---
the droop absorbs reactive power, and a tighter band means it absorbs sooner. In
the importing window the droop must \emph{inject}, and DS voltage instead has an
\emph{interior} minimum at $\dz = \SI{0.0075}{\pu}$
(\SIlist{0.00658;0.00634;0.00587;0.00521;0.00498}{\pu} for
$\dz = 0 \dots 0.0075$, rising thereafter). The narrow dead bands that are
Pareto-optimal under export are therefore \emph{dominated} under import: they
are worse for both objectives at once.

A dead band selected only in the voltage-rise regime would thus have been chosen
against evidence that does not hold when the sign of the interface flow
reverses. Since the import side of the axis reaches only
$-\SI{117}{\mega\watt}$ for the structural reason of
\S\ref{sec:extra-windows}, this conclusion rests on one window and should be
treated as indicative until a deeper-import operating point can be constructed.

\paragraph{The single-objective reading is overturned.} On interface-$Q$
tracking alone the optimum moves $0.005 \rightarrow 0.02 \rightarrow \geq 0.03$
and the wide dead bands appear best. Once DS voltage enters, \emph{every}
$\dz \geq 0.0075$ is dominated in the mildest window: another dead band is
better there in both objectives simultaneously. A dead band that is dominated
in any operating point is a poor fixed choice, because in that condition
voltage quality is surrendered for no tracking benefit.

\paragraph{Compromise point.} Normalising each objective to its own best value
within a window, $r_i(\dz) = f_i(\dz)/\min_\dz f_i(\dz) \geq 1$, and minimising
the worst ratio (Chebyshev scalarisation) gives:

\begin{center}
\begin{tabular}{lcccc}
\toprule
$\dz$ & 2016-01-05 & 2016-01-15 & 2016-12-18 & worst case \\
\midrule
$0$      & 3.16 & 2.87 & 3.17 & 3.17 \\
$0.0025$ & 2.38 & 1.85 & 1.72 & \textbf{2.38} \\
$0.005$  & 2.84 & 1.64 & 2.07 & 2.84 \\
\bottomrule
\end{tabular}
\end{center}

$\dz = \SI{0.0025}{\pu}$ minimises the worst-case relative sacrifice across all
windows and objectives, and is Chebyshev-optimal in two of the three.

\paragraph{Recommendation.} $\dz = \SI{0.005}{\pu}$, and unusually for a
multi-objective problem this needs no weighting argument: it is the only dead
band that is non-dominated in every operating regime tested
(Table~\ref{tab:pareto}). Any other choice is beaten on \emph{both} objectives
simultaneously somewhere in the operating range.

The Chebyshev compromise agrees in two of the four windows and prefers
$\SI{0.001}{\pu}$ in the other two, but that disagreement is not a genuine
alternative: $\SI{0.001}{\pu}$ is dominated under import, so no weighting of the
objectives rescues it there. The scalarisation weights each objective by its own
best value within a window, which is a modelling choice; the Pareto set is free
of that choice, which is why the recommendation is stated from the set rather
than from the scalarisation.

A caveat on one data point. In the importing window the interface-$Q$ error
rises smoothly from \SI{1.334}{\Mvar} at $\dz = 0.005$ to \SI{1.633}{\Mvar} at
$\dz = 0.02$, then drops sharply to \SI{1.178}{\Mvar} at $\dz = 0.03$. That
discontinuity places $0.03$ in the window's Pareto set. It is out of character
with the rest of the curve and should be checked before being relied upon: the
dead zone is known to admit more than one self-consistent equilibrium
(\texttt{docs/qss\_rms\_divergence\_analysis.tex}), so a jump of this kind may
be the run settling into a different one rather than a smooth trend. It does not
affect the recommendation, since $0.03$ is dominated in two other windows.

Produced by \texttt{analysis/deadband\_pareto.py}.

%======================================================================
\section{Disturbance rejection under a load step}
\label{sec:step}

\subsection{Why a second experiment is needed}

The sweep above selects $\dz$ against profile evolution, which is a slow,
continuous excitation. It therefore says nothing about how the dead zone behaves
when the system is displaced abruptly. The two regimes plausibly favour opposite
choices, and that tension is the point of the experiment:

\begin{itemize}
  \item Under \emph{profiled operation} a wider dead band tracked the interface
    better at stressed operating points (\S\ref{sec:findings}), because the
    local droop stops competing with the outer optimisation between dispatches.
  \item Under a \emph{step} a narrower dead band should reject better, because
    the droop begins responding at a smaller voltage deviation and therefore
    acts sooner and more strongly.
\end{itemize}

If both hold, $\dz$ cannot be chosen on tracking quality alone: it also sets how
much of a disturbance the local layer absorbs before the slower cascade
responds.

\subsection{The disturbance is a load step, not a contingency}

The step multiplies every active-load profile
(\texttt{mv\_rural\_pload}, \texttt{HS4\_pload}, \texttt{HS5\_pload}) by a
constant factor from a fixed instant onward. Reactive-load columns are excluded
deliberately: \texttt{mv\_rural\_qload} is signed and predominantly negative in
this data, so scaling it would deepen a capacitive injection rather than add
load.

This is \emph{not} a contingency, and the distinction is structural rather than
terminological. A contingency switches an element out of service by mutating the
pandapower network, which the \rms{} plant rejects --- for that plant the
network is only a measurement mirror, so the change would never reach
PowerFactory and the co-simulation would be silently inconsistent. A load step
needs none of that machinery: it perturbs the exogenous profile, which both
plants already consume through supported paths --- the static plant through
\texttt{apply\_profiles}, the \rms{} plant through
\texttt{Plant.apply\_exogenous} (\texttt{EvtLod} load events). No element is
switched, and both legs see the identical disturbance.

\paragraph{The step is applied after interpolation.} This is an implementation
detail with a substantive consequence. The profile source is
\SI{15}{\minute} data, which \texttt{load\_profiles} interpolates linearly to
the \SI{20}{\second} dispatch step. Perturbing the source file would therefore
produce a \SI{15}{\minute} \emph{ramp}, not a step, and the experiment would
silently measure something else. The perturbation is applied to the interpolated
frame instead. Verified: the ratio between stepped and baseline profiles is
exactly $1$ before the instant and exactly the factor after it, with zero
intermediate samples.

\subsection{Design}

\begin{table}[htbp]
\centering
\begin{tabular}{ll}
\toprule
Item & Value \\
\midrule
Step magnitudes & $\times 1.10$ (mild) and $\times 1.25$ (severe) \\
Step instant & $t = \SI{100}{\second}$ of a \SI{300}{\second} horizon \\
Before / after & 5 / 10 dispatch intervals \\
Dead bands & the full nine-point grid \\
Windows & 2016-01-05 08:00 (mild, sharp interior optimum) and \\
        & 2016-12-18 14:00 (stressed, shallow optimum) \\
Runs & $9 \times 2 \times 2 = 36$ \\
\bottomrule
\end{tabular}
\caption{Disturbance-rejection grid. The \SI{200}{\second} of recovery spans a
full \SI{180}{\second} TSO dispatch period, so the response is not truncated
before the slower layer has acted even once.}
\label{tab:steptable}
\end{table}

The two windows are chosen as the extremes of the profiled-operation result: one
where $\dz^\star$ is sharp and narrow, one where it is shallow and wide. If the
step reverses the ranking, it should be most visible between these two.

\paragraph{Separation from the undisturbed study.} The stepped runs are written
to the same results directory with an otherwise identical configuration. The
admission filter of the profiled-operation analysis therefore requires
\texttt{load\_step\_time\_s} to be absent, so a disturbed run cannot enter the
undisturbed curves. This is the same class of safeguard as the scenario and
multiplier checks of \S\ref{sec:comparability}.

\subsection{Results}

% ===================== RESULTS PLACEHOLDER ============================
% Metrics to report, computed separately over the pre-step and post-step
% intervals rather than averaged across the event:
%   - peak interface-Q excursion after the step
%   - peak nodal-voltage excursion after the step
%   - dispatch intervals to re-settle within the 2% band
%   - steady-state error restored, pre vs post
% ======================================================================

\begin{figure}[htbp]
\centering
% \includegraphics[width=\textwidth]{figures/deadband_step.pdf}
\framebox[0.9\textwidth][c]{\rule{0pt}{6cm}%
  \textsf{PLACEHOLDER: disturbance rejection vs $\dz$}}
\caption{Rejection of the load step against dead-zone half-width, for both step
magnitudes and both windows. Expected: rejection degrades monotonically with
$\dz$, opposing the profiled-operation result.}
\label{fig:step}
\end{figure}

Averaging a metric across the step instant would let the transient dominate the
mean and hide the steady-state behaviour the earlier sections measure, so every
quantity is reported pre- and post-step separately.

%======================================================================
\section{Reporting conventions and threats to validity}

\subsection{Edge values are bounds, not optima}
\label{sec:edge}

An argument minimum that falls on the smallest or largest dead band swept is
not a measured optimum: the metric may still be improving beyond the range. Such
entries are flagged (\texttt{at\_range\_edge}) and reported as bounds. Where the
interface-$Q$ minimum was found at an edge, the range was extended
(\S\ref{sec:windows}) until the minimum was bracketed, or the residual bound is
stated explicitly.

This distinction propagates: if the optima of two windows differ and one is an
edge value, the measured spread is a \emph{lower bound} on the true separation.

\subsection{A flat metric has no optimum}

Where a metric does not vary with $\dz$ within a window, the argument minimum is
an artefact of ordering and is suppressed rather than reported. This is what the
degenerate window of \S\ref{sec:degenerate} would otherwise have contributed.

\subsection{The excitation the dead band was tuned against}
\label{sec:excitation}

The dead zone is selected here against a single class of excitation: the
evolution of the load and DER profiles. No contingency is applied, and
measurement noise is disabled. Two consequences follow, and they act in the
same direction.

\paragraph{Measurement noise.} A dead zone's standard justification is to
prevent the actuator from responding to noise on the measured voltage. That
mechanism is absent here by construction, so the entire narrow-side branch of
the argument --- that a too-small $\dz$ causes the DER to move continuously and
the OFO layer to re-anchor against them --- is carried by genuine profile
variation alone. Since measurement noise would contribute exactly the small,
fast, approximately zero-mean excitation that a dead band exists to reject,
adding it should penalise narrow dead bands further. The value of $\dz^\star$
reported here is therefore expected to be a \emph{lower} bound on the value
appropriate to a noisy plant. The direction of this bias is argued, not
measured; quantifying it requires repeating the sweep with a noise model.

\paragraph{Contingencies.} No element outage is applied anywhere in this study.
The abrupt-disturbance question is addressed instead by the load step of
\S\ref{sec:step}, which reaches both plants through the exogenous-profile path;
a genuine outage would require switching support in the \rms{} plant, which does
not currently exist. A load step and a line trip are not equivalent: the trip
also changes the network topology, and therefore the sensitivities the
controllers hold cached, whereas the step changes only the operating point. The
results here should accordingly not be read as a dead band validated for
post-contingency operation.

\paragraph{Reachability guard.} The quasi-static voltage-reachability guard is
disabled, because it would solve a second plant behind the \rms{} adapter. This
removes a mechanism that could otherwise mask a poor $\dz$ by intervening
before the dead zone's effect became visible in the controlled outputs.

\subsection{Comparability across runs}
\label{sec:comparability}

Runs enter the study only if their stored configuration matches on scenario, DER
capability model, profile use, droop slope, seeding options, and the per-DSO
scenario multipliers. The multiplier check is not cosmetic: an earlier run
predating the $\mathrm{DSO}_3$ multipliers would otherwise have been admitted
beside scaled runs, which the runner itself warns are not comparable.

\subsection{Open items}

\begin{enumerate}
  \item No current-topology severity screening exists
    (\S\ref{sec:representative}); the windows are purposive.
  \item The metric set contains no direct measure of actuator activity, whereas
    the mechanism attributed to a too-narrow dead band is DER chatter and
    repeated re-anchoring. Interface $Q$ registers this only indirectly. A
    direct measure --- DER reactive traverse per dispatch interval, or tap
    operations --- can be derived from the stored per-park records without
    re-running anything, and would evidence the narrow-side branch of the
    argument far better than an inferred penalty.
  \item The sweep has not been repeated with measurement noise enabled
    (\S\ref{sec:excitation}). This is the single change most likely to move
    $\dz^\star$, and in a predictable direction.
  \item Whether the thesis quotes \texttt{rural\_700}, \texttt{base\_410}, or
    both remains open; the results here are \texttt{rural\_700} only.
\end{enumerate}

\end{document}
